% arara: lualatex
% arara: lualatex
\documentclass{article}
\PassOptionsToPackage{colorlinks = true}{hyperref}
\usepackage{
    pline-vec
    ,minted
    ,amssymb
    ,caption
    ,bookmark
    ,hyperref
}
\title{pline-vec}
\author{Jasper}
\date{\today}
\begin{document}
\maketitle
\tableofcontents
% See https://tex.stackexchange.com/a/31787/319072
\section{Introduction}
pline-vec is a 3D toolkit for Ti\textit{k}Z.
\section{Examples}
\subsection{
    Parametric mappings,
    and the pline-vec math libraries
}
\subsubsection{Parametric mappings}
A parametric mapping is a mathematical rule which can accept any 
number of inputs, and produces any number of outputs, all according to
some mathematical formula. For example, A curve in 
\(\mathbb{R}^{3}\) is defined
by a parametric function which accepts one input and delivers three 
outputs. Similarly, a surface in \(\mathbb{R}^{3}\) is defined by 
a parametric function which accepts two inputs, and delivers three
outputs. The parametric equation for a solid accepts three inputs and 
delivers three outputs. Parametric mappings are useful for 
traversing many common mathematical objects. For example,
they can be used to navigate the sphere --- which enables us to 
produce direction vectors in any direction, based on longitude and 
latitude.
\subsubsection{The pline-vec math libraries}
There are two math libraries for pline-vec. The main one is
math.lua, which defines several useful 3D math functions.
In the lua library, to avoid naming conflicts, they are prefixed 
with ``pv\_''. Similarly, they are prefixed with ``pv'' in the pgfmath
library. These are the commands from the lua library:
\begin{itemize}
    \item   
        pv\_cross\_product\_x(u,v)\par
        Delivers the x component of the cross product of 
        two vectors.
    \item   
        pv\_cross\_product\_y(u,v)\par
        Delivers the y component of the cross product of 
        two vectors.
    \item   
        pv\_cross\_product\_z(u,v)\par
        Delivers the z component of the cross product of 
        two vectors.
    \item   
        pv\_dot\_product(u,v)\par
        Delivers the dot product of 
        two vectors.
    \item   
        pv\_norm(u)\par
        Delivers the norm of a vector
    \item 
        pv\_ZYZ\_rotation\_matrix\_x(angles,vector)\par
        Delivers the x component of the ZYZ totation matrix of 
        two vectors.
    \item 
        pv\_ZYZ\_rotation\_matrix\_y(angles,vector)\par
        Delivers the y component of the ZYZ totation matrix of 
        two vectors.
    \item 
        pv\_ZYZ\_rotation\_matrix\_z(angles,vector)\par
        Delivers the z component of the ZYZ totation matrix of 
        two vectors.
    \item 
        pv\_sphere\_x(longitude,latitude)\par
        Delivers the x component of the sphere, based on
        longitude and latitude.
    \item 
        pv\_sphere\_y(longitude,latitude)\par
        Delivers the y component of the sphere, based on
        longitude and latitude.
    \item 
        pv\_sphere\_z(longitude,latitude)\par
        Delivers the z component of the sphere, based on
        longitude and latitude.
\end{itemize}
The pgfmath library is also lua-based, though it is done internally.
\begin{itemize}
    \item 
        pvcrossproductx\par
        x-component of the cross product of u and v,
        \#1 - u1,
        \#2 - u2,
        \#3 - u3,
        \#4 - v1,
        \#5 - v2,
        \#6 - v3
    \item 
        pvcrossproducty\par
        y-component of the cross product of u and v,
        \#1 - u1,
        \#2 - u2,
        \#3 - u3,
        \#4 - v1,
        \#5 - v2,
        \#6 - v3
    \item 
        pvcrossproductz\par
        z-component of the cross product of u and v,
        \#1 - u1,
        \#2 - u2,
        \#3 - u3,
        \#4 - v1,
        \#5 - v2,
        \#6 - v3
    \item 
        pvdotproduct\par
        dot product of u and v,
        \#1 - u1,
        \#2 - u2,
        \#3 - u3,
        \#4 - v1,
        \#5 - v2,
        \#6 - v3
     \item 
        pvnorm\par
        norm product of u
        \#1 - u1,
        \#2 - u2,
        \#3 - u3
    \item 
        pvtransformrotmainx\par
        x-component of ZYZ rotation matrix on a point
        \#1 - alpha,
        \#2 - beta,
        \#3 - gamma,
        \#4 - u1,
        \#5 - u2,
        \#6 - u3
    \item 
        pvtransformrotmainy\par
        y-component of ZYZ rotation matrix on a point
        \#1 - alpha,
        \#2 - beta,
        \#3 - gamma,
        \#4 - u1,
        \#5 - u2,
        \#6 - u3
    \item 
        pvtransformrotmainz\par
        z-component of ZYZ rotation matrix on a point
        \#1 - alpha,
        \#2 - beta,
        \#3 - gamma,
        \#4 - u1,
        \#5 - u2,
        \#6 - u3
    \item 
        pvspherex\par
        x-component of the sphere, given longitude and latitude
        \#1 - longitude,
        \#2 - latitude
    \item 
        pvspherey\par
        y-component of the sphere, given longitude and latitude
        \#1 - longitude,
        \#2 - latitude
    \item 
        pvspherez\par
        z-component of the sphere, given longitude and latitude
        \#1 - longitude,
        \#2 - latitude

\end{itemize}
\subsection{Setting the main orientation}
\subsubsection{Working with sub-orientations}
\subsection{Rendering parametric curves}
\subsection{Rendering parametric surfaces}
\subsection{Rendering parametric solids}
\subsection{
    Rendering parametric curves, 
    surfaces and solids all at once
}
\subsubsection{The problem with tangency}
\subsection{Clipping planes by a rectangular prism}
\subsection{Intersecting planes --- in progress}
\end{document}